\section{Introduction}

Video VLMs spend a surprising amount of time rediscovering what they already
know. A benchmark clip is sampled into several frames, many of those frames are
visually redundant, and the stack still pays again for vision encoding,
prefill, or both. The same waste appears in repeated follow-up questions on the
same video, and it becomes even more obvious in live streams where most of the
scene remains stable for long stretches. The interesting question is no longer
whether redundancy exists. It does. The question is where that redundancy can
be removed without breaking the answer.

That sounds like one problem, but it is really three. First-pass inference is
limited by whatever share of wall-clock the vision tower actually owns.
After-ingest follow-up queries on the same video are a different regime: once
the prefix has been paid, prompt reuse can dominate. Benchmark routing under a
fixed dense backend is different again: it tells us whether a training-free
policy preserves quality, not whether the backend has already been rewritten to
skip wall-clock work. Reuse is not one knob. It is three different regimes
with three different ceilings.

We therefore center a broader anti-recomputation story. The fresh-video
headline is Gemma vision-tower pruning: measured end-to-end gains on VideoMME,
MVBench, and TOMATO, with magnitude governed by a simple
share$\times$reduction ceiling. The largest local multipliers come from
persistent-KV reuse on Qwen, where follow-up questions about the same video
drop to sub-second latency. The routing lane then explains the boundary: why
temporal placement of fresh computation matters more than novelty magnitude,
and why semantic-substitution claims must stay separate from sparse-path
speedup claims.

The same thesis also survives contact with deployment-style workloads. In the
deployment-scale streaming evaluation, same-video follow-up latency is 0.8\,s
on Gemma 4 26B, streaming VideoMME reports 13$\times$ ViT reduction and
\mbox{$\sim$50$\times$} dominant-pipeline reduction, selected real-video
regimes report 4.2--4.5$\times$ end-to-end speedups, and an active-scene RTSP
demo still cuts ViT calls by 14.3$\times$ over a 60\,s window. Those numbers
come from a different protocol and stay labeled that way in the paper. They
matter because they show that the benchmark story survives operational
settings.

We make three contributions:
\begin{itemize}[leftmargin=1.5em]
\item We show that training-free mid-layer vision-tower pruning yields measured
first-pass end-to-end gains on frozen Gemma video-VLM runs, and that those
gains are predicted surprisingly well by a simple ceiling law,
$1 / (1 - V_{\mathrm{share}} V_{\mathrm{red}})$.
\item We show that after-ingest follow-up queries are a much larger reuse
opportunity: persistent-KV reuse on Qwen2.5-VL-7B-4bit reduces same-video
follow-up latency to roughly conversational timescales, with 47.2$\times$
speedup at 8 frames and a clean safe envelope through 16 frames.
\item We provide mechanism and boundary evidence for why these gains are
regime-dependent. On clean-tree Qwen holdouts, a bounded-staleness routing
planner preserves the quality--compute frontier, beats novelty-ranked dense
selection, and shows that fresh-compute placement over time matters more than
novelty magnitude alone.
\end{itemize}

\begin{sectiontodo}{Tighten the opening once the final venue and title settle}
The scientific structure is now the right one, but the opening still needs a
last pass for voice and compression after the final venue template lands. Keep
the dry precision; do not let this turn back into either a Qwen-only methods
note or a second streaming paper.

\textcolor{todoBorder}{Source pointers: \path{paper/priority.md},
\path{paper/claim-matrix.md},
\path{paper/publishability-status.md},
\path{paper/framing.md},
\path{codec-through-sam/paper/publishability-status.md}.}
\end{sectiontodo}
